%% Journal guide: 150-250 words, no undefined abbreviations, no citations. Current count: see wc.
\abstract{Large language models can write interactive 3D simulators of machines from short descriptions. Such programs run and look plausible, yet they can be kinematically wrong, and current pipelines accept them after a vision-language model has looked at screenshots. We present an executable test oracle for generated mechanism simulators. It samples the page's pure rendering function over time and checks structural invariants, the dimensions the simulator declares, engineering design rules, and the kinematic solution computed from those declared dimensions without per-task specifications. Using eight mechanisms, 88 simulators from two models of one vendor (plus 20 from another vendor), we ask which errors occur, what the oracle and a vision-language reviewer detect, whether the oracle can gate iterative repair, and how specific its feedback must be. Under free-form specifications, 20 of 80 simulators were wrong, mostly in design knowledge such as misassigned published link lengths; visual review missed invisible faults such as an impossible gear assembly. In repair of at most three rounds, feedback naming the faulty element with its correct value fixed 28 of 28 faulty cases; naming the element only, 15 of 28; naming the failed check and its magnitude, 11 of 28; self-refinement without feedback, 7 of 28; and the visual reviewer in the loop, 5 of 22. The information needed depends on the repairing model: on design-knowledge errors the smaller model succeeded only when given the correct value (13 of 13 versus 0 of 13), while the larger one needed only the location (7 of 7).}
